\chapter{Parallel Transport}
\label{ch:vii24-parallel-transport}

A connection tells us how to differentiate a vector field along a curve.  Parallel
transport solves the inverse problem: start with one vector and extend it along the
curve so that its covariant derivative vanishes.
\[
\boxed{\frac{DV}{dt}=0.}
\]
For the Levi--Civita connection, this extension preserves lengths and angles.

\section{Parallel vector fields}

\begin{definition}[Parallel field]\label{def:vii24-parallel}
Let \(\gamma:I\to M\) be smooth.  A vector field \(V(t)\in T_{\gamma(t)}M\) is
parallel along \(\gamma\) if
\[
\frac{DV}{dt}
=
\nabla_{\dot\gamma}V
=
0.
\]
\end{definition}

In coordinates,
\[
V=V^k\partial_k,
\]
so the equation becomes
\[
\dot V^k
+
\Gamma^k_{ij}(\gamma(t))\dot x^iV^j
=
0.
\]

\section{Existence and uniqueness}

\begin{theorem}[Parallel initial-value problem]\label{thm:vii24-ivp}
Given a smooth curve \(\gamma:[a,b]\to M\), a time \(t_0\in[a,b]\), and a vector
\(v\in T_{\gamma(t_0)}M\), there exists a unique parallel field \(V\) along
\(\gamma\) satisfying
\[
V(t_0)=v.
\]
\end{theorem}

\begin{proof}
The coordinate equation is a linear first-order ODE for the component vector
\(V(t)\). Standard linear ODE theory gives existence and uniqueness.
\end{proof}

\section{The transport map}

\begin{definition}[Parallel transport]\label{def:vii24-map}
The parallel transport map along \(\gamma\) from \(a\) to \(b\) is
\[
P_\gamma:T_{\gamma(a)}M\to T_{\gamma(b)}M,
\qquad
P_\gamma(v)=V(b),
\]
where \(V\) is the unique parallel field with \(V(a)=v\).
\end{definition}

\begin{proposition}[Linearity and invertibility]\label{prop:vii24-linear}
\(P_\gamma\) is a linear isomorphism.
\end{proposition}

\begin{proof}
Linearity follows from linearity of the parallel ODE.  Reversing the curve solves
the inverse problem, so
\[
P_{\gamma^{-1}}=P_\gamma^{-1}.
\]
\end{proof}

\begin{proposition}[Concatenation]\label{prop:vii24-concat}
If \(\gamma_1\) ends where \(\gamma_2\) begins, then
\[
P_{\gamma_2*\gamma_1}
=
P_{\gamma_2}\circ P_{\gamma_1}.
\]
\end{proposition}

\section{Levi--Civita transport is orthogonal}

\begin{theorem}[Metric preservation]\label{thm:vii24-metric}
For the Levi--Civita connection, if \(V,W\) are parallel along \(\gamma\), then
\[
g(V(t),W(t))
\]
is constant. Consequently
\[
P_\gamma:T_{\gamma(a)}M\to T_{\gamma(b)}M
\]
is an inner-product isometry.
\end{theorem}

\begin{proof}
Metric compatibility gives
\[
\frac d{dt}g(V,W)
=
g(DV/dt,W)+g(V,DW/dt)=0.
\]
\end{proof}

\begin{corollary}[Orthonormal frames remain orthonormal]\label{cor:vii24-frames}
Parallel transport of an orthonormal basis along a curve produces an orthonormal
frame along that curve.
\end{corollary}

\section{Geodesics revisited}

\begin{proposition}[Geodesic tangent]\label{prop:vii24-geodesic}
A curve is an affinely parametrized geodesic exactly when its tangent vector is
parallel along itself:
\[
P_{\gamma|[a,t]}(\dot\gamma(a))=\dot\gamma(t).
\]
\end{proposition}

\section{Transport in a local frame}

Let \(e_1,\dots,e_n\) be a local frame with connection one-form matrix
\(\omega=(\omega^a{}_b)\).  If \(V=V^ae_a\), then parallelism is
\[
\dot V^a
+
\omega^a{}_b(\dot\gamma)V^b
=
0.
\]

\begin{proposition}[Fundamental matrix]\label{prop:vii24-fundamental-matrix}
If \(U(t)\) is the matrix transporting frame components, then
\[
\dot U(t)
=
-A(t)U(t),
\qquad
A^a{}_b(t)=\omega^a{}_b(\dot\gamma(t)),
\qquad
U(a)=I.
\]
Thus \(P_\gamma\) is represented by the endpoint value \(U(b)\).
\end{proposition}

\begin{remark}[Path ordering]\label{rem:vii24-path-order}
If the matrices \(A(t)\) fail to commute at different times, one cannot replace the
solution by an ordinary exponential of \(-\int A\,dt\).  The correct object is the
path-ordered exponential, or equivalently the fundamental solution of the ODE.
\end{remark}

\section{Euclidean examples}

\begin{example}[Cartesian Euclidean transport]\label{ex:vii24-euclidean}
In \(\mathbb R^n\) with Cartesian coordinates, \(\Gamma^k_{ij}=0\). Parallel fields
have constant components, so parallel transport is ordinary translation of
vectors.
\end{example}

\begin{example}[Polar frame warning]\label{ex:vii24-polar}
In Euclidean polar coordinates, a vector with constant \((r,\theta)\)-coordinate
components is generally not parallel because the coordinate frame varies.
Parallelism must be tested with the Christoffel terms.
\end{example}

\section{Parallel transport on a sphere}

Consider the round sphere.  Along a great circle, the tangent vector of the great
circle is parallel. A second parallel vector can be chosen orthogonal to it within
the tangent plane, producing a parallel orthonormal frame along the arc.

\begin{example}[Meridian]\label{ex:vii24-meridian}
Along a meridian on the unit sphere, the unit tangent to the meridian is parallel,
as is the east--west tangent vector obtained by the corresponding Levi--Civita
parallel ODE along that meridian.
\end{example}

\begin{warning}[Transport depends on the path]\label{warn:vii24-path}
On a curved manifold, parallel transport between the same endpoints can depend on
the chosen path.  The systematic study of transport around closed loops is
holonomy, the subject of VII/25.
\end{warning}

\section{Covectors and tensors}

A connection on \(TM\) induces connections on \(T^*M\) and all tensor bundles.

\begin{definition}[Parallel covector]\label{def:vii24-covector}
A covector field \(\alpha(t)\in T^*_{\gamma(t)}M\) is parallel if
\[
\frac{D\alpha}{dt}=0.
\]
Equivalently, for every parallel vector field \(V\),
\[
\alpha(V)
\]
is constant.
\end{definition}

\begin{proposition}[Metric duals]\label{prop:vii24-musical}
For the Levi--Civita connection, \(V\) is parallel if and only if \(V^\flat\) is
parallel.
\end{proposition}

\section{Principal-bundle horizontal lifts}

Let \(P\to M\) be a principal bundle with a principal connection.

\begin{definition}[Horizontal lift of a curve]\label{def:vii24-horizontal-lift}
A curve \(\widetilde\gamma(t)\in P\) above \(\gamma(t)\) is horizontal if
\[
\pi(\widetilde\gamma(t))=\gamma(t),
\qquad
\dot{\widetilde\gamma}(t)\in H_{\widetilde\gamma(t)}P.
\]
\end{definition}

\begin{theorem}[Horizontal lift]\label{thm:vii24-horizontal-lift}
Given \(\gamma\) and an initial point \(p_0\in P_{\gamma(a)}\), there is a unique
horizontal lift \(\widetilde\gamma\) with
\(\widetilde\gamma(a)=p_0\).
\end{theorem}

\begin{remark}
For the frame bundle, horizontal lifting of a frame is precisely parallel
transport of that frame. This unifies the vector-bundle and principal-bundle
pictures developed in VII/08 and VII/21.
\end{remark}

\section{Reparametrization}

\begin{proposition}[Orientation-preserving reparametrization]\label{prop:vii24-reparam}
Parallel transport depends on the oriented image of the curve, not on its regular
orientation-preserving parametrization.
\end{proposition}

\begin{proof}
Under \(t=\phi(s)\), the covariant derivative along the reparametrized curve gains
the scalar factor \(\phi'(s)\), so the zero equation remains zero.
\end{proof}

\section{Exercises}

\begin{exercise}\label{exr:vii24-01}
Define a parallel vector field along a curve.
\end{exercise}
\begin{hint}
% PEDAGOGY-ENRICHED-VII
Set covariant derivative to zero. Make the controlling geometric criterion explicit before the calculation. Solve the parallel-transport equation along the chosen curve and use metric compatibility to check preservation of inner products.
\end{hint}
\begin{solution}
It satisfies \(DV/dt=0\).
\end{solution}


\begin{exercise}\label{exr:vii24-02}
Write the coordinate parallel-transport ODE.
\end{exercise}
\begin{hint}
% PEDAGOGY-ENRICHED-VII
Expand \(V=V^k\partial_k\). Work in a convenient local model first, then return to the invariant statement. Solve the parallel-transport equation along the chosen curve and use metric compatibility to check preservation of inner products.
\end{hint}
\begin{solution}
\(\dot V^k+\Gamma^k_{ij}\dot x^iV^j=0\).
\end{solution}


\begin{exercise}\label{exr:vii24-03}
Why does a unique parallel field exist for prescribed initial data?
\end{exercise}
\begin{hint}
% PEDAGOGY-ENRICHED-VII
Use linear ODE theory. After the main computation, check the hypotheses and geometric interpretation. Solve the parallel-transport equation along the chosen curve and use metric compatibility to check preservation of inner products.
\end{hint}
\begin{solution}
The component equation is a smooth linear first-order system.
\end{solution}


\begin{exercise}\label{exr:vii24-04}
Define the parallel transport map \(P_\gamma\).
\end{exercise}
\begin{hint}
% PEDAGOGY-ENRICHED-VII
Use the endpoint of the unique parallel extension. Make the controlling geometric criterion explicit before the calculation. Solve the parallel-transport equation along the chosen curve and use metric compatibility to check preservation of inner products.
\end{hint}
\begin{solution}
\(P_\gamma(v)=V(b)\) for the parallel field with \(V(a)=v\).
\end{solution}


\begin{exercise}\label{exr:vii24-05}
Why is \(P_\gamma\) linear?
\end{exercise}
\begin{hint}
% PEDAGOGY-ENRICHED-VII
The ODE is linear in \(V\). Work in a convenient local model first, then return to the invariant statement. Solve the parallel-transport equation along the chosen curve and use metric compatibility to check preservation of inner products.
\end{hint}
\begin{solution}
Linear combinations of parallel solutions are parallel with the corresponding initial data.
\end{solution}


\begin{exercise}\label{exr:vii24-06}
Why is \(P_\gamma\) invertible?
\end{exercise}
\begin{hint}
% PEDAGOGY-ENRICHED-VII
Reverse the curve. After the main computation, check the hypotheses and geometric interpretation. Solve the parallel-transport equation along the chosen curve and use metric compatibility to check preservation of inner products.
\end{hint}
\begin{solution}
Transport along the reversed curve is the inverse map.
\end{solution}


\begin{exercise}\label{exr:vii24-07}
State the concatenation law.
\end{exercise}
\begin{hint}
% PEDAGOGY-ENRICHED-VII
Transport in two stages. Make the controlling geometric criterion explicit before the calculation. Solve the parallel-transport equation along the chosen curve and use metric compatibility to check preservation of inner products.
\end{hint}
\begin{solution}
\(P_{\gamma_2*\gamma_1}=P_{\gamma_2}\circ P_{\gamma_1}\).
\end{solution}


\begin{exercise}\label{exr:vii24-08}
Show Levi--Civita transport preserves vector lengths.
\end{exercise}
\begin{hint}
% PEDAGOGY-ENRICHED-VII
Differentiate \(g(V,V)\). Work in a convenient local model first, then return to the invariant statement. Solve the parallel-transport equation along the chosen curve and use metric compatibility to check preservation of inner products.
\end{hint}
\begin{solution}
Parallelism and metric compatibility make the derivative zero.
\end{solution}


\begin{exercise}\label{exr:vii24-09}
Show Levi--Civita transport preserves angles.
\end{exercise}
\begin{hint}
% PEDAGOGY-ENRICHED-VII
Preserve all pairwise inner products. After the main computation, check the hypotheses and geometric interpretation. Solve the parallel-transport equation along the chosen curve and use metric compatibility to check preservation of inner products.
\end{hint}
\begin{solution}
Inner products and norms are constant, hence the angle formula is unchanged.
\end{solution}


\begin{exercise}\label{exr:vii24-10}
What happens to an orthonormal basis under Levi--Civita transport?
\end{exercise}
\begin{hint}
% PEDAGOGY-ENRICHED-VII
Use metric preservation. Make the controlling geometric criterion explicit before the calculation. Solve the parallel-transport equation along the chosen curve and use metric compatibility to check preservation of inner products.
\end{hint}
\begin{solution}
It remains an orthonormal basis along the curve.
\end{solution}


\begin{exercise}\label{exr:vii24-11}
How is a geodesic characterized by parallel transport?
\end{exercise}
\begin{hint}
% PEDAGOGY-ENRICHED-VII
Look at its tangent field. Work in a convenient local model first, then return to the invariant statement. Solve the parallel-transport equation along the chosen curve and use metric compatibility to check preservation of inner products.
\end{hint}
\begin{solution}
Its tangent vector is parallel along the curve.
\end{solution}


\begin{exercise}\label{exr:vii24-12}
Write parallel transport in a local frame.
\end{exercise}
\begin{hint}
% PEDAGOGY-ENRICHED-VII
Use connection one-forms. After the main computation, check the hypotheses and geometric interpretation. Solve the parallel-transport equation along the chosen curve and use metric compatibility to check preservation of inner products.
\end{hint}
\begin{solution}
\(\dot V^a+\omega^a{}_b(\dot\gamma)V^b=0\).
\end{solution}


\begin{exercise}\label{exr:vii24-13}
Write the fundamental-matrix equation.
\end{exercise}
\begin{hint}
% PEDAGOGY-ENRICHED-VII
Collect the linear system. Make the controlling geometric criterion explicit before the calculation. Solve the parallel-transport equation along the chosen curve and use metric compatibility to check preservation of inner products.
\end{hint}
\begin{solution}
\(\dot U=-A(t)U\) with \(A^a{}_b=\omega^a{}_b(\dot\gamma)\) and \(U(a)=I\).
\end{solution}


\begin{exercise}\label{exr:vii24-14}
When can the transport matrix be written as an ordinary exponential?
\end{exercise}
\begin{hint}
% PEDAGOGY-ENRICHED-VII
Require commutation in time. Work in a convenient local model first, then return to the invariant statement. Solve the parallel-transport equation along the chosen curve and use metric compatibility to check preservation of inner products.
\end{hint}
\begin{solution}
If \(A(t_1)\) and \(A(t_2)\) commute for all times, then \(U(b)=\exp(-\int_a^bA(t)\,dt)\).
\end{solution}


\begin{exercise}\label{exr:vii24-15}
Describe Euclidean Cartesian parallel transport.
\end{exercise}
\begin{hint}
% PEDAGOGY-ENRICHED-VII
All Christoffel symbols vanish. After the main computation, check the hypotheses and geometric interpretation. Solve the parallel-transport equation along the chosen curve and use metric compatibility to check preservation of inner products.
\end{hint}
\begin{solution}
Vector components stay constant.
\end{solution}


\begin{exercise}\label{exr:vii24-16}
Why are constant polar-coordinate components not necessarily parallel in the Euclidean plane?
\end{exercise}
\begin{hint}
% PEDAGOGY-ENRICHED-VII
The polar frame varies. Make the controlling geometric criterion explicit before the calculation. Solve the parallel-transport equation along the chosen curve and use metric compatibility to check preservation of inner products.
\end{hint}
\begin{solution}
Nonzero Christoffel symbols contribute to the covariant derivative.
\end{solution}


\begin{exercise}\label{exr:vii24-17}
Along a geodesic, is the tangent field parallel?
\end{exercise}
\begin{hint}
% PEDAGOGY-ENRICHED-VII
Use the geodesic equation. Work in a convenient local model first, then return to the invariant statement. Solve the parallel-transport equation along the chosen curve and use metric compatibility to check preservation of inner products.
\end{hint}
\begin{solution}
Yes: \(D\dot\gamma/dt=0\).
\end{solution}


\begin{exercise}\label{exr:vii24-18}
Define a parallel covector.
\end{exercise}
\begin{hint}
% PEDAGOGY-ENRICHED-VII
Use the induced dual connection. After the main computation, check the hypotheses and geometric interpretation. Solve the parallel-transport equation along the chosen curve and use metric compatibility to check preservation of inner products.
\end{hint}
\begin{solution}
It satisfies \(D\alpha/dt=0\).
\end{solution}


\begin{exercise}\label{exr:vii24-19}
How can parallel covectors be characterized using parallel vectors?
\end{exercise}
\begin{hint}
% PEDAGOGY-ENRICHED-VII
Pair them. Make the controlling geometric criterion explicit before the calculation. Solve the parallel-transport equation along the chosen curve and use metric compatibility to check preservation of inner products.
\end{hint}
\begin{solution}
\(\alpha(V)\) is constant for every parallel \(V\).
\end{solution}


\begin{exercise}\label{exr:vii24-20}
Why does \(\flat\) carry parallel vectors to parallel covectors for Levi--Civita?
\end{exercise}
\begin{hint}
% PEDAGOGY-ENRICHED-VII
Use \(\nabla g=0\). Work in a convenient local model first, then return to the invariant statement. Solve the parallel-transport equation along the chosen curve and use metric compatibility to check preservation of inner products.
\end{hint}
\begin{solution}
Metric compatibility says covariant differentiation commutes with lowering indices.
\end{solution}


\begin{exercise}\label{exr:vii24-21}
Define a horizontal lift in a principal bundle.
\end{exercise}
\begin{hint}
% PEDAGOGY-ENRICHED-VII
Lift the curve with horizontal velocity. After the main computation, check the hypotheses and geometric interpretation. Solve the parallel-transport equation along the chosen curve and use metric compatibility to check preservation of inner products.
\end{hint}
\begin{solution}
It projects to the base curve and has velocity in the horizontal distribution.
\end{solution}


\begin{exercise}\label{exr:vii24-22}
What is the frame-bundle meaning of parallel transport?
\end{exercise}
\begin{hint}
% PEDAGOGY-ENRICHED-VII
Lift a frame horizontally. Make the controlling geometric criterion explicit before the calculation. Solve the parallel-transport equation along the chosen curve and use metric compatibility to check preservation of inner products.
\end{hint}
\begin{solution}
A horizontal frame lift is exactly a frame whose vectors are parallel along the base curve.
\end{solution}


\begin{exercise}\label{exr:vii24-23}
Does parallel transport depend on speed of parametrization?
\end{exercise}
\begin{hint}
% PEDAGOGY-ENRICHED-VII
Use orientation-preserving reparametrization. Work in a convenient local model first, then return to the invariant statement. Solve the parallel-transport equation along the chosen curve and use metric compatibility to check preservation of inner products.
\end{hint}
\begin{solution}
No; regular orientation-preserving reparametrization leaves the zero covariant-derivative equation unchanged.
\end{solution}


\begin{exercise}\label{exr:vii24-24}
State the bridge to holonomy.
\end{exercise}
\begin{hint}
% PEDAGOGY-ENRICHED-VII
Close the curve. After the main computation, check the hypotheses and geometric interpretation. Solve the parallel-transport equation along the chosen curve and use metric compatibility to check preservation of inner products.
\end{hint}
\begin{solution}
VII/25 studies the transformations obtained by parallel transport around loops and the groups they generate.
\end{solution}

\section{Solved problem dossiers}

\begin{problem}[Parallel-ODE problem]\label{prob:vii24-01}
Derive the coordinate ODE for a parallel vector field.
\end{problem}
\begin{solution}
Substitute \(V=V^j\partial_j\) and \(\dot\gamma=\dot x^i\partial_i\) into \(\nabla_{\dot\gamma}V=0\) and apply the connection Leibniz rule.
\end{solution}


\begin{problem}[Existence-uniqueness problem]\label{prob:vii24-02}
Prove the parallel transport map is well defined.
\end{problem}
\begin{solution}
The coordinate equation is a linear ODE with smooth coefficients, so each initial vector has a unique solution along the curve.
\end{solution}


\begin{problem}[Composition problem]\label{prob:vii24-03}
Prove the concatenation law for parallel transport.
\end{problem}
\begin{solution}
Transport the initial vector along the first curve, then use that endpoint as initial data on the second. Uniqueness identifies the piecewise field with transport along the concatenation.
\end{solution}


\begin{problem}[Inverse problem]\label{prob:vii24-04}
Prove transport along the reversed path is the inverse.
\end{problem}
\begin{solution}
Reverse a parallel field in parameter. The reversed field remains parallel, so transporting out and back returns the initial vector.
\end{solution}


\begin{problem}[Metric-preservation problem]\label{prob:vii24-05}
Prove Levi--Civita parallel transport is an isometry.
\end{problem}
\begin{solution}
For two parallel fields \(V,W\), differentiate \(g(V,W)\) and use metric compatibility.
\end{solution}


\begin{problem}[Geodesic problem]\label{prob:vii24-06}
Relate the geodesic equation to parallel transport.
\end{problem}
\begin{solution}
A curve is geodesic precisely when its tangent field solves the parallel ODE along the same curve.
\end{solution}


\begin{problem}[Fundamental-matrix problem]\label{prob:vii24-07}
Represent parallel transport in a moving frame.
\end{problem}
\begin{solution}
The connection matrix evaluated on \(\dot\gamma\) gives \(A(t)\), and the transport matrix solves \(\dot U=-AU\), \(U(a)=I\).
\end{solution}


\begin{problem}[Euclidean polar problem]\label{prob:vii24-08}
Explain how a constant Euclidean vector acquires changing polar components while remaining parallel.
\end{problem}
\begin{solution}
The vector itself is fixed in Cartesian space, but the polar coordinate basis rotates and rescales. Its changing components exactly cancel the Christoffel terms in the parallel ODE.
\end{solution}


\begin{problem}[Sphere problem]\label{prob:vii24-09}
Construct a parallel orthonormal frame along a great-circle arc.
\end{problem}
\begin{solution}
The unit tangent is parallel because the arc is geodesic. Choose an orthogonal unit tangent vector at the start and parallel transport it; metric preservation keeps the frame orthonormal.
\end{solution}


\begin{problem}[Dual-transport problem]\label{prob:vii24-10}
Describe transport on covectors.
\end{problem}
\begin{solution}
Define \(P_\gamma^*\) contragrediently so pairings with transported vectors remain constant; this agrees with parallel transport for the induced dual connection.
\end{solution}


\begin{problem}[Horizontal-lift problem]\label{prob:vii24-11}
Explain how a principal connection produces transport.
\end{problem}
\begin{solution}
Lift the base curve horizontally from a chosen point in the initial fiber. The endpoint defines an equivariant map between fibers, and associated-bundle representations turn it into vector transport.
\end{solution}


\begin{problem}[Path-dependence problem]\label{prob:vii24-12}
Why can two paths with the same endpoints produce different transport maps?
\end{problem}
\begin{solution}
The connection prescribes infinitesimal comparison, but curvature can make the accumulated comparison path dependent. Closed-loop discrepancies are organized by holonomy.
\end{solution}

\section{Challenge problems}

\begin{problem}[Path-ordered exponential]\label{prob:vii24-ch01}
Explain why noncommuting connection matrices require path ordering.
\end{problem}
\begin{solution}
For \(\dot U=-A(t)U\), successive infinitesimal factors occur in chronological order. If \(A(t_1)A(t_2)\neq A(t_2)A(t_1)\), rearranging them changes the product, so the solution is a time/path-ordered exponential rather than \(\exp(-\int A)\).
\end{solution}


\begin{problem}[Gauge covariance of transport]\label{prob:vii24-ch02}
How does the matrix of parallel transport change under a frame change \(e'=eA\)?
\end{problem}
\begin{solution}
If endpoint frames change by \(A(a)\) and \(A(b)\), the transport matrix changes by \(U'\!=A(b)^{-1}UA(a)\). Thus the underlying linear map is frame independent.
\end{solution}


\begin{problem}[Flat connection and path independence]\label{prob:vii24-ch03}
On a simply connected domain with a globally trivial flat connection that admits a parallel frame, why is transport path independent?
\end{problem}
\begin{solution}
In a global parallel frame the connection matrix vanishes, so parallel components are constant. The transport map then depends only on the endpoint frame identification, not on the chosen path.
\end{solution}


\begin{problem}[Parallel tensors]\label{prob:vii24-ch04}
Show that tensor products and contractions of parallel tensor fields remain parallel.
\end{problem}
\begin{solution}
The induced connection obeys a Leibniz rule on tensor products and commutes with contractions. If each factor has zero covariant derivative, so does the tensor product, and contraction preserves zero.
\end{solution}


\begin{problem}[Holonomy preview]\label{prob:vii24-ch05}
Why does transport around loops based at \(p\) naturally form a subgroup of the orthogonal group \(O(T_pM,g_p)\)?
\end{problem}
\begin{solution}
Each loop gives a Levi--Civita transport isometry \(T_pM\to T_pM\). Concatenation gives composition, reversal gives inverses, and the constant loop gives the identity. VII/25 studies this group in detail.
\end{solution}

\section*{Bridge to VII/25}
Parallel transport around closed loops need not return a vector unchanged. VII/25 organizes these loop transports into holonomy groups and relates them to curvature and global geometry.
